\chapter{Énigmes sur les nombres}\label{ch:iv1:number-riddles}

\begin{ivtrains}
Problèmes sur les nombres entiers~: les aiguilles d'une horloge, les factorielles, les restes, les derniers chiffres, la divisibilité et les montants que l'on peut former avec des pièces données. Ils récompensent le travail modulo un petit nombre et la recherche d'une structure plutôt que l'étude de nombreux cas. Les premières questions sont des échauffements, et les dernières demandent une vraie persévérance.
\end{ivtrains}

\begin{question}[$\star$ \iqroles{QT, Trader, Analyste} \iqfirm{toutes} \iqlevel{junior}]\label{iq:iv1:number-riddles:1}
Quel est l'angle entre l'aiguille des heures et celle des minutes d'une horloge à 4~h~36~? Les deux aiguilles se déplacent en continu.
\end{question}

\begin{question}[$\star$ \iqroles{QR, QT, Trader, QD, Dev} \iqfirm{toutes} \iqlevel{junior}]\label{iq:iv1:number-riddles:2}
Combien de zéros termine $250!$ lorsqu'il est écrit en chiffres décimaux ordinaires~?
\end{question}

\begin{question}[$\star$ \iqroles{QR, QT, Trader} \iqfirm{toutes} \iqlevel{junior}]\label{iq:iv1:number-riddles:3}
Sans calculatrice, quel est le dernier chiffre de $7^{2026}$ lorsqu'il est écrit en décimal~?
\end{question}

\begin{question}[$\star$ \iqroles{QR, QT, Trader, QD, Dev} \iqfirm{toutes} \iqlevel{junior}]\label{iq:iv1:number-riddles:4}
Un couloir compte 1000 portes fermées, numérotées de 1 à 1000. Pour chaque $k$ de 1 à 1000, une personne passe et change l'état de chaque porte dont le numéro est un multiple de $k$. Combien de portes sont ouvertes à la fin~? Combien de fois la porte 72 est-elle changée~?
\end{question}

\begin{question}[$\star$ \iqroles{QR, QT, Trader} \iqfirm{toutes} \iqlevel{junior}]\label{iq:iv1:number-riddles:5}
Prenez un nombre quelconque à trois chiffres et écrivez-le deux fois de suite pour obtenir un nombre à six chiffres, comme 123123. Pourquoi le résultat est-il toujours divisible par 7, 11 et 13~?
\end{question}

\begin{question}[$\star$ \iqroles{QR, QT, Trader} \iqfirm{toutes} \iqlevel{junior}]\label{iq:iv1:number-riddles:6}
Quel est le plus petit entier naturel non nul qui laisse le reste 1 dans la division par 2, le reste 2 dans la division par 3, et ainsi de suite jusqu'au reste 9 dans la division par 10~?
\end{question}

\begin{question}[$\star\star$ \iqroles{QR, QT, Trader, QD} \iqfirm{toutes} \iqlevel{junior}]\label{iq:iv1:number-riddles:7}
Combien de zéros termine $20!$ lorsqu'il est écrit en base 12~?
\end{question}

\begin{question}[$\star\star$ \iqroles{QR, QT, Trader} \iqfirm{toutes} \iqlevel{junior}]\label{iq:iv1:number-riddles:8}
Le nombre à cinq chiffres $4a27b$, où $a$ et $b$ sont des chiffres, est divisible par 72. Trouvez $a$ et $b$.
\end{question}

\begin{question}[$\star\star$ \iqroles{QR, QT, QD} \iqfirm{toutes} \iqlevel{junior}]\label{iq:iv1:number-riddles:9}
Quel est le reste de la division de $3^{2026}$ par 17~? Faites le calcul à la main.
\end{question}

\begin{question}[$\star\star$ \iqroles{QR, QT, Trader} \iqfirm{toutes} \iqlevel{junior}]\label{iq:iv1:number-riddles:10}
Un nombre laisse le reste 1 dans la division par chacun des nombres 2, 3, 4, 5 et 6, et il est divisible par 7. Quel est le plus petit entier naturel non nul de ce type, et quel est le suivant~?
\end{question}

\begin{question}[$\star\star$ \iqroles{QR, QT, Trader, QD} \iqfirm{teneur de marché, hedge fund} \iqlevel{junior}]\label{iq:iv1:number-riddles:11}
Vous disposez de pièces de 7 et de pièces de 11, en nombre illimité. Quel est le plus grand montant que vous ne pouvez pas payer exactement~? Combien de montants strictement positifs ne pouvez-vous pas payer~?
\end{question}

\begin{question}[$\star\star$ \iqroles{QR, QT, Trader, QD, Dev} \iqfirm{toutes} \iqlevel{junior}]\label{iq:iv1:number-riddles:12}
Vous écrivez tous les entiers de 1 à 1000. Combien de fois écrivez-vous le chiffre 1~? Quelle est la somme de tous les chiffres des nombres de 1 à 999~?
\end{question}

\begin{question}[$\star\star$ \iqroles{QR, QT, Trader, QD, Dev} \iqfirm{toutes} \iqlevel{junior}]\label{iq:iv1:number-riddles:13}
Quel est le plus petit jeu de poids, chacun étant un nombre entier de grammes, qui permet à une balance de peser tout objet dont la masse est un nombre entier de grammes de 1 à 40, si un poids peut être posé sur l'un ou l'autre plateau~? Combien de poids faut-il s'ils ne peuvent être posés que sur un seul plateau~?
\end{question}

\begin{question}[$\star\star$ \iqroles{QR, QT, QD} \iqfirm{toutes} \iqlevel{junior}]\label{iq:iv1:number-riddles:14}
Quels sont les deux derniers chiffres de $3^{1000}$~? Faites le calcul à la main.
\end{question}

\begin{question}[$\star\star$ \iqroles{QR, QT} \iqfirm{teneur de marché, hedge fund} \iqlevel{junior}]\label{iq:iv1:number-riddles:15}
Sans calculer le nombre, quelle est la plus grande puissance de 2 qui divise $\binom{100}{50}$~? Quelle est la plus grande puissance de 3~? $\binom{100}{50}$ est-il divisible par 7~?
\end{question}

\begin{question}[$\star\star$ \iqroles{QR, QT} \iqfirm{toutes} \iqlevel{junior}]\label{iq:iv1:number-riddles:16}
Montrez que, pour tout nombre premier $p$ strictement supérieur à 3, $p^2 - 1$ est divisible par 24.
\end{question}

\begin{question}[$\star\star\star$ \iqroles{QR, QT, Trader, QD} \iqfirm{teneur de marché, hedge fund} \iqlevel{junior, confirmé}]\label{iq:iv1:number-riddles:17}
Trouvez le nombre à quatre chiffres qui, multiplié par 4, donne le même nombre lorsque ses chiffres sont lus dans l'ordre inverse.
\end{question}

\begin{question}[$\star\star\star$ \iqroles{QR, QT, Trader} \iqfirm{teneur de marché, hedge fund} \iqlevel{junior, confirmé}]\label{iq:iv1:number-riddles:18}
Considérez chaque décimal périodique $0.abcabcabc\dots$ où $a$, $b$ et $c$ sont trois chiffres distincts, le zéro étant autorisé. Quelle est la somme de tous ces nombres~?
\end{question}

\begin{question}[$\star\star\star$ \iqroles{QR, QT, Trader} \iqfirm{teneur de marché, hedge fund} \iqlevel{junior, confirmé}]\label{iq:iv1:number-riddles:19}
Une horloge a des aiguilles des heures et des minutes de même longueur, de sorte qu'on ne peut pas les distinguer, et les deux se déplacent en continu. Pendant un cycle de 12 heures, à combien de moments ne peut-on pas lire l'heure d'après la position des aiguilles~?
\end{question}

\begin{question}[$\star\star\star$ \iqroles{QR, QT, Trader, QD} \iqfirm{teneur de marché, hedge fund} \iqlevel{junior, confirmé}]\label{iq:iv1:number-riddles:20}
Des boîtes contiennent 6, 9 ou 20 articles, et vous devez acheter des boîtes entières. Quel est le plus grand nombre d'articles que vous ne pouvez pas acheter exactement~? Montrez que tout nombre plus grand peut être acheté.
\end{question}

\begin{question}[$\star\star\star$ \iqroles{QR, QT} \iqfirm{teneur de marché, hedge fund} \iqlevel{junior, confirmé}]\label{iq:iv1:number-riddles:21}
Quels entiers naturels non nuls ne peuvent pas s'écrire comme une somme d'au moins deux entiers naturels non nuls consécutifs~? Démontrez-le. De combien de façons peut-on écrire 2026 comme une telle somme~?
\end{question}

\begin{question}[$\star\star\star$ \iqroles{QR, QT, QD, Dev} \iqfirm{teneur de marché, hedge fund} \iqlevel{junior, confirmé}]\label{iq:iv1:number-riddles:22}
100 personnes se tiennent en cercle, numérotées de 1 à 100 dans l'ordre. En faisant le tour du cercle à partir de la personne 1, la personne 1 reste, la personne 2 sort, la personne 3 reste, la personne 4 sort, et ainsi de suite, en sautant toujours une personne restante et en éliminant la suivante, jusqu'à ce qu'il ne reste qu'une personne. Quel numéro reste-t-il~?
\end{question}

\begin{question}[$\star\star\star\star$ \iqroles{QR, QT} \iqfirm{teneur de marché, hedge fund} \iqlevel{confirmé}]\label{iq:iv1:number-riddles:23}
Quel est le plus petit entier naturel non nul qui double lorsque son dernier chiffre est déplacé en tête~? Par exemple, déplacer le dernier chiffre de 1234 en tête donne 4123. Le nombre ne peut pas commencer par 0.
\end{question}

\begin{question}[$\star\star\star\star$ \iqroles{QR, QT} \iqfirm{teneur de marché, hedge fund} \iqlevel{confirmé}]\label{iq:iv1:number-riddles:24}
Soit $A$ la somme des chiffres décimaux de $4444^{4444}$. Soit $B$ la somme des chiffres de $A$, et soit $C$ la somme des chiffres de $B$. Que vaut $C$~? Trouvez-le à la main.
\end{question}

\begin{question}[$\star\star\star\star$ \iqroles{QR, QT, Trader} \iqfirm{teneur de marché, hedge fund} \iqlevel{confirmé}]\label{iq:iv1:number-riddles:25}
Cinq marins réunissent un tas de noix de coco. Pendant la nuit, chacun se réveille une fois, à tour de rôle. Chacun partage le tas en cinq parts égales avec une noix de coco en trop, donne cette noix à un singe, cache une part et rassemble les quatre autres en un tas. Le matin, les noix restantes se partagent exactement en cinq parts égales. Quel est le plus petit nombre possible de noix de coco dans le tas initial~?
\end{question}
